\chapter{Conclusion}%
\label{ch:conclusion}

% The Conclusion chapter is typically the last chapter you write, though most often it will be the first chapter others will read (after the abstract).
% Make sure to be concise and complete, and make it not too detailed.
% No new information is presented or derived, all information should have been presented in other chapters.

% \section{Summary \& Findings}
% Here you summarise your work (the context, questions, approach), and present the findings in context with the research questions.

% \section{Contributions}
% Highlight the contributions your thesis provides.
% These are tangible results, which can live outside the context of this thesis.
% Examples are new methods, frameworks, software.
% Make sure to provide links to Git repos or datasets if applicable.

% \section{Limitations \& Threats to Validity}
% No work is perfect.
% Discuss the limitations of your study.
% A research study can only be as unbiased as the researcher and the circumstances they are working with.
% The threats to validity offers insights into the variables that were encountered and how you compensated for them.

% \section{Future Work}
% This is the moment to highlight all the work you wanted to do if you had more time.
% Or any new ideas you did not have the time to explore.
% This is also where you can give an idea how to address any limitations mentioned in the previous section.